\documentclass[11pt]{article}
\usepackage{amsmath,amssymb,amsthm}
\usepackage[left=2cm,top=2cm,right=2cm,bottom=2cm,head=2cm,foot=1cm]{geometry}
\title{Variant Classification}

\newcommand{\p}{~\\}
\newcommand{\Sm}{S_{m\bold{\cdot}}}

\begin{document}
\maketitle

\begin{center}
  \begin{tabular}{c|l}
    $Y_i$     & Phenotype of individual $i$ ($0=$ unaffected, $1=$ affected)\\
    $G_{il}$   & Genotype at individual $i$ locus $l$\\
    $\phi_l$  & Phenotypic effect of non-reference allele at locus $l$\\
  \end{tabular}
\end{center}

\section{Introduction}

\p In variant classification, we use phenotype and genotype
observations to decide which mutations should be reported as
pathogenic. Example scenarios are:

\paragraph{Scenario 1: A VUS associated with disease incidence in relatives:}
We have performed exon sequencing of the relevant disease genes for an
individual and discovered a Variant of Uncertain Significance
(VUS). Initially we lack any further evidence regarding whether or not
the variant is causative. Subsequently, we discover the same VUS in
other individuals, some of which have affected relatives. How many
other individuals, with what sorts of patterns of family disease
incidence, are sufficient to conclude that the variant should be
reported as likely deleterious?

\paragraph{Scenario 2: Follow-up sequencing of an affected relative:}
We have performed exon sequencing of the relevant disease genes for an
individual and discovered a VUS. The individual has two living
relatives with diagnoses of the disease. We sequence the relatives and
discover the same VUS in one of them. Is this sufficient evidence to
report the variant as likely deleterious? (Note that sequencing a
single affected individual is equivalent to sequencing individual A
who has affected relative B, with the coefficient of relatedness
between A and B equal to 1.)

\p Consider a VUS at locus $l$ which has been observed in several
individuals. A variant classification algorithm should have the
following properties:

\begin{enumerate}

\item The strength of evidence for pathogenicity of $l$ increases with
  the number of families in which there are affected individuals
  related to an individual with the VUS

\item The evidence for pathogenicity of $l$ is stronger if the
  affected individuals are more closely related to the sequenced
  individual with the VUS.

\end{enumerate}


\section{Model and variant classification algorithm}

\p Let $\phi_l$ represent the phenotypic effect of the non-reference
genotypes at locus $l$. We can initially consider an extremely simple
disease model: a locus $l$ has either $\phi_l=0$ (non-causative) or
$\phi_l=1$ (causative); any individual with 1 or more non-reference
alleles in any causative locus will certainly have the disease
phenotype from birth. (dominant, early onset, fully
penetrant). Otherwise the individual will certainly not display the
disease phenotype. (The sequenced genes explain all disease
incidence.)

\p The variant classification problem is inference for the effect
indicators $\phi_l$. The pedigree in each family is known. We
introduce latent ``meiosis indicator'' variables: $S_{ml}=0$ if
meiosis $m$ transmitted the maternal allele at locus $l$; $S_{ml}=1$
if the paternal allele was transmitted. $\Sm$ denotes the sequence of
meiosis indicator values for all loci at meiosis $m$.

\p The likelihood is

\begin{align*}
\Pr(Y, G|\phi)
=& \sum_S \Pr(S) \Pr(Y,G|S,\phi)\\
=& \sum_S \Pr(S) \Pr(G|S) \Pr(Y|G,\phi)\\
=& \sum_S \Big(\prod_m\Sm \Big) \Pr(G|S) \Pr(Y|G,\phi)\\
\end{align*}

\p In words: we compute the likelihood by averaging over all possible
histories of recombination breakpoints in the pedigrees, weighted by
their prior probabilities (these can be computed easily from the
recombination map). Given a particular history $S$ of meiosis
indicators, the contribution to the likelihood is the product of the
prior probability $\Pr(S)$ of that meiotic history, the probability
$\Pr(G|S)$ of the observed genotypes given the meiotic indicators, and
the probability $\Pr(Y|G,\phi)$ of the phenotype data given the
genotypes.


\end{document}
